\documentclass[11pt]{article}

% ----------------------------------------------------------------------
%  A BEGINNER'S GUIDE to Robust Workload Feasibility in CSLAP
%  Plain-language companion to this repository + the
%  iscf10kt CPLEX-vs-HiGHS replication.  Self-contained; compile on Overleaf.
% ----------------------------------------------------------------------

\usepackage[utf8]{inputenc}
\usepackage[T1]{fontenc}
\usepackage{amsmath,amssymb,amsthm}
\usepackage{booktabs}
\usepackage{geometry}
\usepackage{enumitem}
\usepackage{xcolor}
\usepackage[hidelinks]{hyperref}
\geometry{a4paper,margin=2.4cm}

% Theorem-like environments
\theoremstyle{definition}
\newtheorem{definition}{Definition}
\theoremstyle{plain}
\newtheorem{proposition}{Proposition}
\theoremstyle{remark}
\newtheorem{remark}{Remark}

% A simple "plain words" box that needs no extra package
\newcommand{\plain}[1]{%
  \par\medskip\noindent\fbox{\parbox{0.97\linewidth}{%
  \textbf{In plain words.}\; #1}}\par\medskip}

\newcommand{\Lnom}{L^{0}}   % nominal pick-line count
\newcommand{\Lhat}{\hat{L}} % deviation magnitude

\title{\bf Robust Workload Feasibility in Warehouse\\ Storage Assignment (CSLAP)\\[4pt]
\large A Beginner's Guide to the Theory, the Model, and the Results}
\author{Companion notes for the CSLAP research project}
\date{\today}

\begin{document}
\maketitle

\begin{abstract}
This note explains, from scratch, a piece of warehouse-optimization research.
The warehouse problem is CSLAP: decide where to store each product so that
orders are picked with as few station visits as possible, without overloading
any station. The research asks a robustness question: \emph{if the demand we
trained on changes later, does a good storage plan stay feasible, and can we
protect it in advance?} We define every term, write out every equation, explain
the robust-optimization theory (Bertsimas--Sim ``budget of uncertainty''), and
then walk through all the experimental results in plain language. No prior
knowledge of optimization is assumed beyond high-school algebra and the idea of
``minimize something subject to rules''.
\end{abstract}

\tableofcontents

% ======================================================================
\section{The setting: what problem are we solving?}
% ======================================================================

Imagine an automated ``pick-and-pass'' warehouse. Products sit in slots at a
number of \emph{picking stations}. A customer \emph{order} is a small basket of
products. To fulfil an order, a tote travels past the stations; it must
\emph{stop} at every station that holds at least one of the ordered products.
Each stop is a \emph{visit}, and visits cost time and money.

Two intuitions drive everything:
\begin{itemize}[nosep]
  \item If two products are frequently ordered \emph{together}, storing them at
  the \emph{same} station saves a visit. So a good layout keeps
  ``correlated'' products together.
  \item Each station has limits: a maximum number of slots, and a maximum amount
  of picking \emph{work} it can do per planning period.
\end{itemize}

The \textbf{Correlated Storage Location Assignment Problem (CSLAP)} is: assign
every product to exactly one station so that the total number of station visits
over all orders is as small as possible, while respecting the slot and workload
limits.

\plain{Put popular-together products on the same station to cut visits, but do
not stuff any one station beyond its slot count or its work capacity.}

% ======================================================================
\section{Notation (the symbols we will use)}
% ======================================================================

\begin{center}
\begin{tabular}{ll}
\toprule
Symbol & Meaning \\
\midrule
$P$ & set of products (SKUs); a single product is $p$ \\
$S$ & set of picking stations; a single station is $s$ \\
$O$ & set of orders; a single order is $o$ \\
$P_o \subseteq P$ & the products contained in order $o$ (its ``support'') \\
$\zeta_s$ & slot capacity of station $s$ (how many products fit) \\
$V_s$ & processing speed of station $s$ (lines per unit time) \\
$T_s$ & time (workload) capacity of station $s$ \\
$L_p$ & \emph{pick-line count} of product $p$: number of distinct orders that contain $p$ \\
$x_{ps}$ & decision: $1$ if product $p$ is stored at station $s$, else $0$ \\
$y_{os}$ & auxiliary: $1$ if order $o$ visits station $s$, else $0$ \\
$\Phi_s$ & the set of products actually placed at station $s$, i.e. $\{p : x_{ps}=1\}$ \\
\bottomrule
\end{tabular}
\end{center}

\begin{definition}[Pick line]
A \emph{pick line} for product $p$ is one order that asks for $p$. So $L_p$
counts how many orders want product $p$. A ``heavy'' product has a large $L_p$.
\end{definition}

\begin{definition}[Workload of a station]
The workload of station $s$ under a layout is the total pick-lines of the
products stored there, divided by the station speed:
$W_s = \sum_{p \in \Phi_s} L_p / V_s$. It measures how much picking work that
station must do.
\end{definition}

% ======================================================================
\section{The nominal placement model (the exact optimization problem)}
% ======================================================================

We now write CSLAP as a \emph{mixed-integer linear program} (MILP): a set of
linear equations and inequalities with some yes/no ($0/1$) variables, plus one
quantity to minimize.

\begin{align}
\min_{x,\,y}\quad & \sum_{o \in O}\sum_{s \in S} y_{os} \label{eq:obj}\\
\text{s.t.}\quad
& \sum_{s \in S} x_{ps} = 1 && \forall p \in P, \label{eq:assign}\\
& \sum_{p \in P} x_{ps} \le \zeta_s && \forall s \in S, \label{eq:slot}\\
& x_{ps} \le y_{os} && \forall o \in O,\; p \in P_o,\; s \in S, \label{eq:link}\\
& \sum_{p \in P} \frac{L_p}{V_s}\, x_{ps} \le T_s && \forall s \in S, \label{eq:work}\\
& x_{ps}\in\{0,1\},\quad y_{os}\in[0,1] && \forall p,o,s. \label{eq:dom}
\end{align}

Let us read each line.

\begin{itemize}
  \item \textbf{Objective \eqref{eq:obj}} counts total visits: for every order
  $o$ and station $s$, $y_{os}=1$ adds one visit. Minimizing the sum minimizes
  visits.
  \item \textbf{Assignment \eqref{eq:assign}} says each product is stored at
  exactly one station (not zero, not two).
  \item \textbf{Slot capacity \eqref{eq:slot}} says a station holds at most
  $\zeta_s$ products.
  \item \textbf{Visit linking \eqref{eq:link}} says: if product $p$ of order $o$
  sits at station $s$ ($x_{ps}=1$), then order $o$ must visit $s$
  ($y_{os}\ge 1$). Because the objective \emph{pushes $y$ down}, at the optimum
  $y_{os}$ equals $1$ exactly when \emph{some} product of $o$ is at $s$ --- this
  is a logical OR, and no ``big-$M$'' trick is needed.
  \item \textbf{Workload \eqref{eq:work}} keeps each station's work under its
  capacity $T_s$.
  \item \textbf{Domain \eqref{eq:dom}}: $x$ is yes/no; $y$ may be a fraction but
  automatically comes out $0/1$ at the optimum, so we can leave it continuous
  (this makes the model much easier to solve).
\end{itemize}

\plain{Equation \eqref{eq:obj} is the goal (fewest visits). Equations
\eqref{eq:assign}--\eqref{eq:work} are the rules: place each product once, do not
exceed slots, count a visit whenever a station holds an ordered product, and do
not overload a station's work capacity.}

\subsection{Where does the workload capacity $T_s$ come from?}

The project sets $T_s$ from the \emph{training} demand with a $10\%$ safety
margin, spread evenly over the stations:
\begin{equation}
T_s = \left\lceil \sigma \cdot \frac{\sum_{p\in P} L_p^{\mathrm{tr}}}{V_s\,|S|}\right\rceil,
\qquad \sigma = 1.10. \label{eq:ceiling}
\end{equation}
Here $L_p^{\mathrm{tr}}$ is the pick-line count measured on the training window,
$|S|$ is the number of stations, and $\lceil\cdot\rceil$ rounds up. In words:
take the total training work, divide it fairly among stations, and allow each
one $10\%$ slack.

\plain{$T_s$ is a ``fair share plus $10\%$'' work budget, computed once from the
data we planned on.}

% ======================================================================
\section{The research question: demand drift and workload feasibility}
% ======================================================================

We build the layout on \emph{training} demand. Later, the warehouse runs on
\emph{new} demand. \emph{Demand drift} means the popularity of products changes:
some products get heavier ($L_p$ grows), others lighter. A layout that respected
\eqref{eq:work} on the training data might now overload a station.

The question of this study is:
\begin{quote}\itshape
Does a layout that satisfies the (training) workload constraint remain
workload-feasible on later demand? If not, how big are the deviations, and how
much ``protection'' would we have needed to stay safe?
\end{quote}

This is different from the project's earlier hypothesis (called H1), which tried
to make the \emph{visit} objective robust and was falsified. Here we protect the
\emph{workload constraint} instead. The objective is left untouched.

% ======================================================================
\section{Part A --- An audit: the ``clean'' feasibility check can be an illusion}
% ======================================================================

Before proposing any fix, the study audits how the project \emph{measures}
out-of-sample feasibility. This is a subtle but important point.

The evaluation compares the \emph{test-window} workload of each station against
the \emph{training-calibrated} ceiling $T_s$ from \eqref{eq:ceiling}, and counts
how many stations exceed it (\texttt{wl\_broken}).

\subsection{The volume artifact}

The temporal experiment uses \emph{expanding} training windows: train on the
first $50\%$ of orders and test on the last $50\%$; then $60/40$, $70/30$,
$80/20$, $90/10$. As the split moves later:
\begin{itemize}[nosep]
  \item the training window grows, so $T_s$ (built from training volume)
  \emph{grows};
  \item the test window shrinks, so the test workload \emph{shrinks}.
\end{itemize}
We are then comparing a tiny test workload against a ceiling meant for a huge
training volume. The ceiling becomes far too loose, and \texttt{wl\_broken}
$=0$ ``by construction'', regardless of how unbalanced the layout is.

The study captures this with a \emph{bindingness} number:
\begin{equation}
B \;=\; \frac{\sum_p L_p^{\mathrm{te}}}{\sum_s V_s\,T_s}
\;=\; \frac{\text{total test volume}}{\text{total ceiling capacity}}. \label{eq:binding}
\end{equation}
If $B < 1/\sigma \approx 0.909$, the check essentially cannot fire on a balanced
layout. In the data $B$ collapses from about $1.0$ down to $0.11$ across the
cuts, so late-cut ``feasible'' verdicts are just \emph{volume slack}, not
balance.

\plain{On late cuts there is so little test demand that no station can look
overloaded --- ``feasible'' there means ``not much happened'', not ``well
balanced''.}

\subsection{The fair fix: a volume-normalized ceiling}

To compare fairly, re-scale the ceiling to the \emph{test} volume:
\begin{equation}
T^{\mathrm{n}}_s = \left\lceil \sigma \cdot \frac{\sum_{p} L_p^{\mathrm{te}}}{V_s\,|S|}\right\rceil.
\label{eq:tnorm}
\end{equation}
This asks: \emph{if the test window carried a normal amount of demand, would a
station be overloaded?} Under this fair ceiling the picture reverses: on the
$480$-product instance every temporal cut has $2$--$4$ of $8$ stations over the
line, with peak loads up to $2.5\times$ the fair share.

\plain{Judged fairly (same volume basis), the trained layouts are \emph{not}
balanced under drift --- some stations would be badly overloaded if demand
returned to normal. This ``fragility'' is what the robust method targets. Note
it is a \emph{what-if}: it is about latent fragility, not overload that
literally happened in a low-volume test window.}

% ======================================================================
\section{Part B --- A robust-optimization primer}
% ======================================================================

\emph{Robust optimization} means: instead of trusting a single guess for
uncertain numbers, we plan for a whole \emph{set} of possible values and require
the plan to work for every value in that set.

The uncertain numbers here are the pick-line counts $L_p$: tomorrow's demand may
differ from today's estimate. Write the estimate as $\Lnom_p$ (the training
value) and let $\Lhat_p \ge 0$ be a ``how much could it move'' amount. The
uncertain value lives in an interval:
\begin{equation}
\tilde L_p \in \bigl[\,\Lnom_p - \Lhat_p,\; \Lnom_p + \Lhat_p\,\bigr]. \label{eq:interval}
\end{equation}

If we protect against \emph{every} product hitting its worst value at once, we
get the extremely cautious \emph{box} (Soyster) model --- usually too
conservative, because it is very unlikely that \emph{all} products spike
simultaneously.

\subsection{The Bertsimas--Sim ``budget of uncertainty''}

Bertsimas and Sim proposed a dial. Choose a whole number $\Gamma$ (the
\emph{budget}) and assume that \emph{at most $\Gamma$ products deviate at the
same time}. Then:
\begin{itemize}[nosep]
  \item $\Gamma = 0$: nobody deviates $\Rightarrow$ the ordinary model.
  \item $\Gamma$ large: everybody may deviate $\Rightarrow$ the full box.
  \item $\Gamma$ in between: a tunable amount of caution.
\end{itemize}

Because the workload constraint is a ``$\le$'' with nonnegative coefficients,
only \emph{upward} moves ($\Lnom_p + \Lhat_p$) can hurt us. So the worst case at
station $s$ adds the sum of the $\Gamma$ \emph{largest} deviations among the
products stored there. The robust workload constraint is:
\begin{equation}
\sum_{p\in\Phi_s}\frac{\Lnom_p}{V_s}\,x_{ps}
\;+\;
\underbrace{\max_{\substack{J\subseteq\Phi_s\\ |J|\le\Gamma}}\ \sum_{p\in J}\frac{\Lhat_p}{V_s}}_{\text{worst-case extra load}}
\;\le\; T_s
\qquad\forall s\in S. \label{eq:robustraw}
\end{equation}

\plain{``Keep the normal load \emph{plus} the biggest $\Gamma$ demand spikes
under capacity.'' $\Gamma$ decides how many simultaneous spikes you insure
against.}

\subsection{Turning the worst case into linear algebra}

The inner ``pick the $\Gamma$ largest'' is a small optimization by itself. We
can write it as a linear program (LP): give each stored product a weight
$z_p\in[0,1]$, choose weights summing to at most $\Gamma$, and maximize the
weighted deviation:
\begin{equation}
\psi_s(x;\Gamma) \;=\; \max_{z}\ \sum_{p} \frac{\Lhat_p}{V_s}\,x_{ps}\,z_p
\quad\text{s.t.}\quad \sum_p z_p \le \Gamma,\ \ 0\le z_p\le 1. \label{eq:innerLP}
\end{equation}
Because all coefficients are nonnegative, the best choice puts $z_p=1$ on the
$\Gamma$ heaviest deviations --- exactly the ``sum of the $\Gamma$ largest''.
(When $\Gamma$ is a whole number this LP gives the same value as the discrete
problem, so nothing is lost.)

% ======================================================================
\section{The exact linear robust counterpart (Proposition 1)}
% ======================================================================

The constraint \eqref{eq:robustraw} contains a ``max inside'', which solvers do
not like. Linear-programming \emph{duality} lets us replace that inner maximum
with a few ordinary linear inequalities. This is the heart of the method.

\begin{proposition}[Exact linear counterpart]\label{prop:1}
Fix a station $s$, a whole-number budget $\Gamma\ge 0$, and an assignment
$x_{\cdot s}\in\{0,1\}^{P}$. The robust constraint \eqref{eq:robustraw} holds
\emph{if and only if} there exist numbers $\theta_s\ge 0$ and $\mu_{ps}\ge 0$
(for every product $p$) such that
\begin{align}
\sum_{p\in P}\frac{\Lnom_p}{V_s}\,x_{ps} \;+\; \Gamma\,\theta_s \;+\; \sum_{p\in P}\mu_{ps} &\;\le\; T_s, \label{eq:dual1}\\
\theta_s + \mu_{ps} &\;\ge\; \frac{\Lhat_p}{V_s}\,x_{ps} \qquad \forall p\in P. \label{eq:dual2}
\end{align}
Since $\Lhat_p/V_s$ and $\Lnom_p/V_s$ are data constants, both lines are
\emph{linear} in the variables $(x,\theta,\mu)$.
\end{proposition}

\begin{proof}[Why it is true (sketch)]
The inner LP \eqref{eq:innerLP} has a \emph{dual} LP. Give a price $\theta_s$ to
the budget row ``$\sum z_p \le \Gamma$'' and a price $\mu_{ps}$ to each cap
``$z_p \le 1$''. The dual says: choose the cheapest prices so that, for every
product, $\theta_s + \mu_{ps}$ covers that product's deviation
$\tfrac{\Lhat_p}{V_s}x_{ps}$; its cost is $\Gamma\theta_s + \sum_p \mu_{ps}$.
\emph{Strong duality} guarantees this cheapest cost \emph{equals} the worst-case
extra load $\psi_s(x;\Gamma)$. Substituting that equality into
\eqref{eq:robustraw} gives exactly \eqref{eq:dual1}--\eqref{eq:dual2}.
\end{proof}

\plain{We replace ``the sum of the $\Gamma$ biggest spikes'' by a small pricing
scheme: a price $\theta_s$ for the budget and a slack $\mu_{ps}$ per product.
Duality proves the cheapest prices reproduce the true worst case \emph{exactly}
--- so the robust rule becomes ordinary linear inequalities a solver can
handle.}

\begin{remark}[Why the ``at most $\Gamma$ of the stored products'' is not a problem]
The set $\Phi_s$ of stored products is itself a decision. This normally makes
robust models hard (``endogenous uncertainty''). Here it causes no trouble:
because the factor $x_{ps}$ multiplies $\Lhat_p$ inside \eqref{eq:innerLP}, a
product that is \emph{not} stored contributes $0$, so we may let the sums range
over all products without changing anything. The dualization then applies
verbatim.
\end{remark}

% ======================================================================
\section{The full robust placement model}
% ======================================================================

Applying Proposition~\ref{prop:1} at every station and substituting for
\eqref{eq:work} gives the complete robust MILP:
\begin{align}
\min_{x,y,\theta,\mu}\quad & \sum_{o\in O}\sum_{s\in S} y_{os} \label{eq:R-obj}\\
\text{s.t.}\quad
& \sum_{s} x_{ps}=1 && \forall p, \label{eq:R-assign}\\
& \sum_{p} x_{ps}\le \zeta_s && \forall s, \label{eq:R-slot}\\
& x_{ps}\le y_{os} && \forall o,\,p\in P_o,\,s, \label{eq:R-link}\\
& \sum_{p}\frac{\Lnom_p}{V_s}x_{ps} + \Gamma_s\,\theta_s + \sum_{p}\mu_{ps} \le T_s && \forall s, \label{eq:R-work}\\
& \theta_s + \mu_{ps} \ge \frac{\Lhat_p}{V_s}x_{ps} && \forall p,s, \label{eq:R-dual}\\
& x_{ps}\in\{0,1\},\ y_{os}\in[0,1],\ \theta_s\ge0,\ \mu_{ps}\ge0. \label{eq:R-dom}
\end{align}
Only the workload block changed: lines \eqref{eq:R-obj}--\eqref{eq:R-link} are
identical to the nominal model. The cost of robustness is one small
\emph{continuous} block ($\theta_s,\mu_{ps}$), not new yes/no variables --- so
the problem does not get combinatorially harder.

\begin{proposition}[Properties of the dial]\label{prop:2}
For whole-number $\Gamma$:
\begin{enumerate}[label=(\roman*),nosep]
  \item $\Gamma=0$ recovers the nominal model \eqref{eq:obj}--\eqref{eq:work}.
  \item If $\Gamma\ge\zeta_s$ for every $s$, the constraint becomes the full box
  (Soyster): $\sum_p \tfrac{\Lnom_p+\Lhat_p}{V_s}x_{ps}\le T_s$.
  \item Increasing $\Gamma$ only \emph{removes} feasible layouts, so the best
  achievable visit count can only rise. More protection is never cheaper.
\end{enumerate}
\end{proposition}

\plain{Turning the dial $\Gamma$ up from $0$ smoothly moves from ``ignore
uncertainty'' to ``guard against everything at once'', and always costs (weakly)
more visits.}

% ======================================================================
\section{How large are the deviations? Choosing $\Lhat_p$ and $\Gamma$}
% ======================================================================

The model needs the deviation sizes $\Lhat_p$. This is the hardest, most honest
part: we must estimate them \emph{before} seeing the future, using only training
data. The study tries three recipes:
\begin{equation}
\text{M1: } \Lhat_p = c_1\,\Lnom_p,\qquad
\text{M2: } \Lhat_p = c_2\sqrt{\Lnom_p},\qquad
\text{M3: } \Lhat_p = \operatorname{std}_{b=1,\dots,B}\bigl(L_p^{(b)}\bigr).
\label{eq:models}
\end{equation}
\begin{itemize}[nosep]
  \item \textbf{M1} assumes deviations grow in proportion to demand
  (multiplicative drift).
  \item \textbf{M2} assumes ``counting noise'' ($\sqrt{\cdot}$ scaling, as for
  Poisson counts).
  \item \textbf{M3} splits the training window into $B$ time-blocks and takes the
  per-product standard deviation of the block counts. \emph{M3 is the only recipe
  a planner can actually compute ahead of time} --- M1/M2 constants are fit to
  realized deviations and are used only as diagnostics.
\end{itemize}

\subsection{Three quantities you need first: \texorpdfstring{$\bar W_s$}{Wbar}, \texorpdfstring{$W^{\mathrm n}_s$}{Wn}, and \texorpdfstring{$n$}{n}}

The two diagnostics below (and most of the results) use three quantities. All
three are ordinary warehouse counts.

\begin{description}
  \item[\quad Nominal load $\bar W_s$ ---``the work we planned for''.]
  Take the products placed at station $s$, add up their \emph{training}
  pick-lines, divide by the speed: $\bar W_s = \sum_{p\in\Phi_s}\Lnom_p/V_s$.
  This is the workload the layout was \emph{designed} around. Example: if station
  $7$ holds $60$ products whose training pick-lines total $41{,}000$ (and
  $V_s=1$), then $\bar W_7 = 41{,}000$ ``grab'' operations over the planning
  window.

  \item[\quad Realized normalized load $W^{\mathrm n}_s$ ---``the work that
  actually showed up''.] The same sum, but with the \emph{future} (test)
  pick-lines $L^{\mathrm n}_p$, rescaled to the same total volume so the
  comparison is apples-to-apples: $W^{\mathrm n}_s=\sum_{p\in\Phi_s}L^{\mathrm
  n}_p/V_s$. If the products at station $7$ became more popular, then
  $W^{\mathrm n}_7 > \bar W_7$ --- say $W^{\mathrm n}_7 = 77{,}000$. That gap,
  $77{,}000$ vs.\ the planned $41{,}000$, \emph{is} the demand drift, felt at one
  station.

  \item[\quad Station occupancy $n$ ---``how many products live here''.]
  $n=|\Phi_s|$ is simply the number of products stored at station $s$, at most
  $\zeta_s$. If all $60$ slots are used, $n=60$. It appears in some formulas
  because a station's workload is a \emph{sum of $n$ product demands}, so how much
  that sum can randomly wobble depends on $n$ (more on this in
  Remark~\ref{rem:prob}).
\end{description}

\paragraph{A running warehouse example (reused throughout).}
Eight stations, $60$ slots each (all full, so $n=60$), equal speed $V_s=1$.
Training demand totals $8\times 41{,}000$ pick-lines, so the fair share is
$41{,}000$ per station and the ceiling is $T_s=\lceil 1.10\times
41{,}000\rceil\approx 45{,}000$. A balanced training layout puts each station at
$\bar W_s\approx 41{,}000$, comfortably under $45{,}000$. Now the future demand
\emph{drifts}: a handful of products at \textbf{station 7} become much more
popular, pushing its realized load to $W^{\mathrm n}_7 = 77{,}000$. Suppose the
per-product increases (``deviations'') at station 7, largest first, are
\[
\Lhat\text{-realized}: \quad +14{,}000,\ +10{,}000,\ +7{,}000,\ +4{,}000,\ +2{,}000,\ \text{then small}.
\]
We will read the diagnostics off this one station.

\subsection{The two retrospective diagnostics: cover and bind}

Both ask a ``what budget would we have needed?'' question, \emph{after} seeing the
drift (hence \emph{retrospective}). Recall $\psi_s(x;\Gamma)$ from
\eqref{eq:innerLP} is just ``the sum of the $\Gamma$ largest deviations at station
$s$'' --- the extra load the budget-$\Gamma$ protection reserves room for.

\medskip
\noindent\textbf{$\Gamma_{\mathrm{cover}}(s)$ --- how many spikes to \emph{see} the overload.}
\begin{equation}
\Gamma_{\mathrm{cover}}(s) = \min\Bigl\{\Gamma :\ \underbrace{\bar W_s}_{\text{planned}} + \underbrace{\psi_s(x;\Gamma)}_{\text{$\Gamma$ biggest spikes}} \ \ge\ \underbrace{W^{\mathrm n}_s}_{\text{what happened}}\Bigr\}.
\end{equation}
It is the smallest budget whose reserved room reaches the load that actually
occurred --- i.e.\ the number of product-spikes we would have had to insure for
the robust constraint to have \emph{predicted} this overload.

\emph{Worked out at station 7:} we need planned $41{,}000$ plus spikes to reach
$77{,}000$, i.e.\ $+36{,}000$ of deviation.
\[
\Gamma{=}1:\,+14{,}000\ (55{,}000);\quad
\Gamma{=}2:\,+24{,}000\ (65{,}000);\quad
\Gamma{=}3:\,+31{,}000;\quad
\Gamma{=}4:\,+35{,}000;\quad
\Gamma{=}5:\,+37{,}000\ (78{,}000)\ \ge 77{,}000.
\]
So $\Gamma_{\mathrm{cover}}(7)=5$: insuring the \emph{five} biggest movers (out
of $60$ products) would have covered this overload. A \emph{small}
$\Gamma_{\mathrm{cover}}$ means the overload was caused by a few big movers, not
by diffuse noise --- exactly the situation the budget model is good at.

\medskip
\noindent\textbf{$\Gamma_{\mathrm{bind}}(s)$ --- how much protection before the layout must change.}
\begin{equation}
\Gamma_{\mathrm{bind}}(s) = \min\{\Gamma :\ \bar W_s + \psi_s(x;\Gamma) > T_s\}.
\end{equation}
It is the smallest budget at which the protected left-hand side breaks the
\emph{training} ceiling $T_s$ --- i.e.\ the point where adding protection would
force the optimizer to \emph{rebalance the layout} at training time.

\emph{Worked out at station 7:} planned $41{,}000$ plus spikes must exceed
$T_s=45{,}000$, i.e.\ just $+4{,}000$. But the single largest spike is already
$+14{,}000$, so $\Gamma{=}1$ already gives $55{,}000 > 45{,}000$. Hence
$\Gamma_{\mathrm{bind}}(7)=1$: even insuring \emph{one} product blows the training
ceiling here, so \emph{any} nonzero budget forces the layout to move products
around. A \emph{small} $\Gamma_{\mathrm{bind}}$ means protection ``bites''
immediately; a large one means the station has slack to absorb spikes for free.

\plain{$\Gamma_{\mathrm{cover}}$ = ``how many product-spikes would we have had to
insure to catch this overload'' (small = a few culprits). $\Gamma_{\mathrm{bind}}$
= ``how little insurance already forces us to rebuild the layout'' (small =
protection is expensive here). In the example, $\Gamma_{\mathrm{cover}}=5$ and
$\Gamma_{\mathrm{bind}}=1$: cheap to detect, but it forces a rebuild at once.}

\begin{remark}[A probabilistic reading --- why a small budget is usually enough]
\label{rem:prob}
Why not just protect against \emph{all} $n$ products spiking at once (the paranoid
box)? Because that essentially never happens. Picture each of the $n$ products at
a station as a coin: on any given future day its demand lands a bit \emph{above}
or a bit \emph{below} its estimate, roughly at random. For the station to
overload, an unusually large number of coins must come up ``high''
\emph{simultaneously}. Just as flipping $60$ fair coins almost never gives $50$
heads, many products spiking together is very unlikely --- the ups and downs
mostly cancel. Formally, if the deviations are independent and symmetric,
\begin{equation}
\Pr\Bigl[\underbrace{\textstyle\sum_{p\in\Phi_s}\tilde L_p/V_s}_{\text{station's actual work}} > T_s\Bigr]
\;\lesssim\; \exp\!\bigl(-\Gamma^2/(2n)\bigr),\qquad n=|\Phi_s|,
\end{equation}
so the risk of overload \emph{shrinks very fast} as the budget $\Gamma$ grows.
Because random ups and downs partly cancel, the \emph{typical} net swing of a sum
of $n$ items grows like $\sqrt n$, not like $n$ (the ``random-walk'' rule). So a
budget of about $\sqrt n$ already covers the realistic worst case.

\emph{Example.} With $n=60$ products at a station, $\sqrt{60}\approx 8$. Insuring
the $8$ biggest possible spikes ($\Gamma=8$) already makes an overload very
unlikely --- yet it is far cheaper than the paranoid box that guards all $60$.
That ``$\Gamma\sim\sqrt n$ buys most of the safety'' is the whole appeal of the
Bertsimas--Sim dial. (Real pick-line counts are not perfectly symmetric, so this
is only a guide; the study leans on the data-driven diagnostics above instead.)
\end{remark}

% ======================================================================
\section{How the model is actually solved}
% ======================================================================

Even though the model is ``only'' linear, the stations are identical, which
creates \emph{symmetry} that stalls exact branch-and-bound. So the solver
pipeline is:
\begin{enumerate}[label=(\arabic*),nosep]
  \item a \textbf{feasibility-first greedy} start (place heaviest products first,
  then swap to repair any workload violation);
  \item a \textbf{deterministic local search} that swaps products to reduce
  visits, keeping the robust constraint as a hard filter;
  \item a \textbf{MILP solver} (HiGHS or CPLEX) warm-started with that
  incumbent, used mainly for bounds and for \emph{proofs of infeasibility}.
\end{enumerate}
The visit objective is measured on the \emph{real} orders; the training supports
are only used to shape the placement. Two solver backends are available and give
identical models --- this lets us check that conclusions are not a solver
artifact (Section~\ref{sec:iscf10kt}).

% ======================================================================
\section{Results}
% ======================================================================

We report four things: the audit (Part A), the robust experiment on a small
drifting instance, a stationary ``control'' benchmark, and a second drifting
instance solved with two different solvers.

\subsection{How to read every number (with the running example)}
\label{sec:metrics}

Before any table, here is what each column means, in warehouse terms, using the
same eight-station example (fair ceiling $T^{\mathrm n}=45{,}000$; after drift,
station 7 carries $77{,}000$, and say two other stations creep to $48{,}000$ and
$46{,}000$, the rest sit at or below $45{,}000$).

\begin{description}
  \item[\quad Bindingness $B$ (Eq.\,\ref{eq:binding}) --- ``is the test even big
  enough to judge?''] The ratio of total test work to total ceiling capacity.
  $B\approx 1$ means the test window carries about as much work as the ceilings
  assume, so an overload \emph{can} show up. $B\ll 1$ means the test window is
  tiny (few orders), so \emph{no} station can look full --- any ``feasible''
  verdict there is meaningless. It is a trust meter for the whole check, not a
  property of the layout.

  \item[\quad Peak $=\displaystyle\max_s \frac{W^{\mathrm n}_s}{T^{\mathrm n}_s}$
  --- ``how bad is the worst station?''] Take the busiest station's realized
  workload and divide by the fair ceiling. The ceiling already includes the
  $10\%$ slack, so:
  \begin{itemize}[nosep]
    \item peak $=1.0$ $\Rightarrow$ the worst station is \emph{exactly} at its
    allowed limit;
    \item peak $=1.71$ $\Rightarrow$ it is carrying $71\%$ \emph{more} work than
    allowed (our station 7: $77{,}000/45{,}000=1.71$);
    \item peak $=2.5$ $\Rightarrow$ two-and-a-half times its fair ceiling --- a
    serious bottleneck.
  \end{itemize}
  Lower is better; anything above $1.0$ is an overloaded station. Peak is a
  \emph{severity} measure (how bad the worst case is), not a count.

  \item[\quad Violation count (``viol.\ /8'') --- ``how many stations are over
  the line?''] Simply the number of stations whose realized load exceeds the
  ceiling. In our snapshot, three stations ($77{,}000$, $48{,}000$, $46{,}000$)
  are above $45{,}000$, so the count is $3$ out of $8$.
  \emph{Why do tables show a fraction like $2.6$?} Because each result is an
  \emph{average over the five temporal cuts}. If the five cuts had counts
  $3,3,3,2,2$, the mean is $2.6$ --- you cannot have $2.6$ overloaded stations in
  a single warehouse; it just means ``about two-to-three of the eight, typically''.
  Violation count is an \emph{incidence} measure (how many), complementary to
  peak (how bad).

  \item[\quad Excess fraction $X$ --- ``how much of the work spills over the
  ceilings?''] Add up, over all stations, the workload sitting \emph{above} the
  ceiling, and divide by the total workload:
  $X=\sum_s (W^{\mathrm n}_s-T^{\mathrm n}_s)^+ / \sum_s W^{\mathrm n}_s$. In the
  snapshot the overflow is $(77{-}45)+(48{-}45)+(46{-}45)=36$ (thousands) out of
  $\approx 360$ total, so $X\approx 10\%$. It answers ``what share of all the
  picking work is piled above where stations should cap out'' --- a volume view,
  between peak (one station) and count (how many).

  \item[\quad PoR (Price of Robustness) --- ``the insurance premium, paid in
  visits''] The percent \emph{increase} in \emph{training} visits caused by
  switching on the robust budget, versus the plain $\Gamma=0$ layout. Example
  from the real data: the nominal layout needs $222{,}270$ visits on the training
  orders; the $\Gamma=1$ robust layout needs $224{,}453$, so
  $\mathrm{PoR}=(224{,}453-222{,}270)/222{,}270 = +0.98\%$. You pay about $1\%$
  more stops to buy the protection.

  \item[\quad $G$ (out-of-sample gain) --- ``did protection also help on real
  future orders?''] The percent \emph{decrease} in \emph{test} visits versus
  $\Gamma=0$. Positive $G$ = the robust layout also picks future orders with fewer
  visits; negative $G$ = it costs \emph{more} visits out of sample too. Real
  example: nominal needs $237{,}287$ test visits, robust needs $238{,}328$, so
  $G=(237{,}287-238{,}328)/237{,}287=-0.44\%$. \emph{$G$ is negative throughout
  this study}: the robust budget buys workload safety, never cheaper visits.
\end{description}

\plain{Three ways to look at overload: \textbf{peak} = how bad is the worst
station; \textbf{violation count} = how many stations are over (a fraction like
$2.6$ is just the average across the five cuts); \textbf{$X$} = what share of all
the work spills over. And two prices: \textbf{PoR} = extra visits you pay to
train the robust layout; \textbf{$G$} = whether it also helps on real future
orders (here it never does --- it insures feasibility, not visits).}

% ----------------------------------------------------------------------
\subsection{The audit, on the 480-product instance}
% ----------------------------------------------------------------------

Table~\ref{tab:audit} shows the two ways of checking feasibility side by side,
over the five expanding temporal cuts (metrics defined in
Section~\ref{sec:metrics}). Read the columns as:
\begin{itemize}[nosep]
  \item \textbf{Cut} --- which train/test split. Cut 0 is $50/50$ (train on the
  first half of orders, test on the second half); cut 4 is $90/10$ (train on
  $90\%$, test on the last $10\%$). Later cuts have a \emph{smaller} test window.
  \item \textbf{$B$} --- bindingness: is the test window big enough for the check
  to mean anything? ($B$ near $1$ = yes; $B$ near $0$ = the test is too small to
  overload anything.)
  \item \textbf{as-run viol.} --- stations flagged using the \emph{training}
  ceiling (the project's original check).
  \item \textbf{norm.\ viol.} --- stations flagged using the \emph{fair}
  (test-volume) ceiling $T^{\mathrm n}_s$ of Eq.\,\eqref{eq:tnorm}.
  \item \textbf{norm.\ peak} and \textbf{$X$} --- severity and spill-over under
  the fair ceiling.
\end{itemize}

\begin{table}[h]\centering
\caption{iscf480, five temporal cuts. As-run vs. volume-normalized checks.}
\label{tab:audit}
\begin{tabular}{cccccc}
\toprule
Cut & $B$ & as-run viol.\ (/8) & norm.\ viol.\ (/8) & norm.\ peak & $X$ (\%) \\
\midrule
0 & 1.005 & 4 & 3 & 1.71 & 11.0 \\
1 & 0.693 & 2 & 3 & 1.91 & 19.8 \\
2 & 0.457 & 1 & 3 & 2.25 & 26.6 \\
3 & 0.281 & 0 & 2 & 2.52 & 27.2 \\
4 & 0.125 & 0 & 3 & 2.38 & 30.7 \\
\bottomrule
\end{tabular}
\end{table}

\emph{Reading two rows.} \textbf{Cut 0} ($50/50$): $B=1.005$, so the test window
is as heavy as the ceilings assume and the check is meaningful --- both views
roughly agree ($4$ vs $3$ stations over). \textbf{Cut 4} ($90/10$): $B=0.125$, so
the test window carries only about an eighth of the assumed work; the training
ceiling is now hopelessly loose, and the as-run check reports $0$ stations over
--- ``perfectly clean''. But fairly measured (fair ceiling), $3$ of $8$ stations
are still overloaded, the worst at $2.38\times$ its ceiling, with $30.7\%$ of the
(little) test work spilling over.

Now look down the two violation columns. The \textbf{as-run} column marches
$4,2,1,0,0$ --- it looks like feasibility keeps \emph{improving} as we move to
later cuts. But that is an illusion: it is tracking $B$ falling
$1.005\to 0.125$, i.e.\ the test window shrinking, \emph{not} the layout getting
better. The \textbf{fair} column stays at $2$--$3$ overloaded stations while the
peak actually \emph{worsens} ($1.71\to 2.52$). Same layouts, opposite stories.

\plain{The ``all clean'' late cuts are a mirage created by a shrinking test
window ($B\to0$). Judged on a fair, equal-volume basis, $2$--$3$ of the $8$
stations are overloaded on \emph{every} cut, by up to $2.5\times$ --- and the
picture gets worse, not better, exactly where the original check says
``feasible''.}

% ----------------------------------------------------------------------
\subsection{The robust experiment, on the 480-product instance}
% ----------------------------------------------------------------------

Table~\ref{tab:iscf480} compares arms (means over the cuts). Here $\alpha$ is a
partial-insurance factor: the model is built with $\alpha\,\Lhat_p$ instead of
$\Lhat_p$, because full-size protection does not fit the $10\%$ slack.

\begin{table}[h]\centering
\caption{iscf480 robust arms (means over the five cuts).}
\label{tab:iscf480}
\begin{tabular}{lccccc}
\toprule
Arm & PoR (\%) & $G$ (\%) & viol.\ (/8) & peak & $X$ (\%) \\
\midrule
$\Gamma=0$ (nominal) & 0.00 & 0.00 & 2.6 & 1.77 & 17.4 \\
$\Gamma=1,\ \alpha=0.5$ & 2.34 & $-4.44$ & 2.6 & \textbf{1.59} & \textbf{13.0} \\
tighten $\beta=1.02$ (same price) & 2.26 & $-1.73$ & 2.8 & 1.75 & 16.9 \\
Hexaly reference & $-3.88$ & $6.16$ & 2.8 & 2.15 & 23.1 \\
\bottomrule
\end{tabular}
\end{table}

Findings:
\begin{itemize}[nosep]
  \item \textbf{Cheap partial insurance helps severity.} $\Gamma=1$ at half
  strength cuts the peak $1.77\to1.59$ and $X$ from $17.4\%\to13.0\%$ for a
  $2.3\%$ visit price --- and the peak falls on all five cuts.
  \item \textbf{Blunt tightening does not.} Lowering the ceiling for everyone at
  the \emph{same} price leaves the peak at $1.75$. The budget helps because it
  specifically avoids co-locating two big ``spiky'' products; uniform tightening
  only rescales the average.
  \item \textbf{Full protection does not fit.} At full strength the solver
  \emph{proves} infeasibility for $\Gamma\ge4$: the drift is simply bigger than
  the $10\%$ headroom. Insuring it fully would require buying more capacity.
  \item \textbf{No free visits.} Every robust arm has $G<0$: it costs visits out
  of sample. The thing being insured is workload feasibility, not visits.
  \item \textbf{Better visit optimizers concentrate load.} The strong Hexaly
  layout wins visits but has the \emph{worst} peak ($2.15$): squeezing visits
  packs correlated, drifting products together.
\end{itemize}

% ----------------------------------------------------------------------
\subsection{A control: the stationary article benchmark (exp02a)}
% ----------------------------------------------------------------------

To isolate \emph{why} protection sometimes helps, the same solver was run on 29
synthetic benchmark instances (50--2000 products) that are \emph{near-stationary}
(little drift). Table~\ref{tab:exp02a} summarizes.

\begin{table}[h]\centering
\caption{exp02a benchmark (means; $\beta=1.02$ is uniform tightening).}
\label{tab:exp02a}
\begin{tabular}{lccccc}
\toprule
Size & arm & PoR (\%) & viol.\ /$|S|$ & peak (max) & infeasible \\
\midrule
50   & $\Gamma=0$ & 0.00 & 0.67/5  & 1.00 (1.05) & 0 \\
50   & $\Gamma=4$ & 2.29 & 0.00    & 0.95 (0.98) & 0 \\
500  & $\Gamma=4$ & 0.13 & 0.20    & 0.99 (1.00) & 0 \\
1000 & $\Gamma=4$ & 0.94 & 0.50    & 1.00 (1.00) & 0 \\
2000 & $\Gamma=4$ & 0.47 & 0.33    & 1.00 (1.02) & 0 \\
\bottomrule
\end{tabular}
\end{table}

Here the nominal layout is already essentially feasible out of sample
(violations are hairline). Full protection is \emph{free} ($\text{PoR}<1\%$, no
infeasibility) --- the \emph{opposite} of the drifting instance, where $\Gamma\ge
4$ was impossible.

\plain{On stationary data there is nothing to insure, so protection costs
nothing and gains little. The value \emph{and} the cost of the robust workload
constraint are both driven by demand drift. This is a clean ``negative
control'': it confirms \emph{why} the method sometimes helps, without claiming it
always does.}

% ----------------------------------------------------------------------
\subsection{A second drifting instance and two solvers (iscf10kt)}
\label{sec:iscf10kt}
% ----------------------------------------------------------------------

The strongest positive claim above was ``targeted budget beats equal-price
uniform tightening''. A single instance is weak evidence, so we re-ran the
experiment on a larger drifting instance (4328 products, 8 stations, five
temporal cuts), using the honest, leakage-free calibration (M3 at full strength,
$\alpha=1$), \emph{and} solved it with \emph{both} the HiGHS and the CPLEX
backends.

\paragraph{The two solvers agree exactly.}
Across all 25 arms, CPLEX and HiGHS returned \emph{identical} numbers (maximum
difference $0.0$ on every metric). At this scale both solvers simply return the
greedy+local-search warm start (they cannot close the large optimality gap in
the time budget). So the solver choice does not change any conclusion --- but it
also means the numbers rest on a metaheuristic incumbent, which is about $21.6\%$
worse on visits than the strong Hexaly layout.

\paragraph{The dominance claim does \emph{not} replicate.}
Table~\ref{tab:iscf10kt} shows, per cut, the peak of the targeted budget
($\Gamma=1$) versus equal-price uniform tightening ($\beta=1.02$).

\begin{table}[h]\centering
\caption{iscf10kt: targeted budget vs.\ uniform tightening, out-of-sample peak
(identical under CPLEX and HiGHS). ``winner'' = lower peak.}
\label{tab:iscf10kt}
\begin{tabular}{ccccc}
\toprule
Cut & nominal peak & $\Gamma=1$ peak (PoR) & $\beta=1.02$ peak (PoR) & winner \\
\midrule
0 & 1.381 & \textbf{1.158} ($-1.1\%$) & 1.255 ($+1.0\%$) & $\Gamma=1$ \\
1 & 1.488 & 1.404 ($-0.5\%$) & \textbf{1.346} ($+1.7\%$) & tighten \\
2 & 1.495 & 1.426 ($+1.4\%$) & \textbf{1.405} ($+1.2\%$) & tighten \\
3 & 1.893 & \textbf{1.534} ($-0.0\%$) & 1.658 ($-0.2\%$) & $\Gamma=1$ \\
4 & 2.129 & 2.015 ($+1.4\%$) & \textbf{1.965} ($+0.1\%$) & tighten \\
\bottomrule
\end{tabular}
\end{table}

\begin{itemize}[nosep]
  \item \textbf{Mechanism survives:} $\Gamma=1$ lowers the peak versus nominal on
  \emph{all five} cuts (e.g.\ $1.893\to1.534$).
  \item \textbf{Competitive claim fails:} $\Gamma=1$ beats uniform tightening on
  peak in only $2$ of $5$ cuts; tightening wins or ties in the other $3$. On the
  small instance the budget won $5/5$; here it is a coin-flip.
  \item \textbf{Full protection still does not fit:} $\Gamma=2$ was
  undetermined/infeasible on $4$ of $5$ cuts. Where it did fit (cut 0) it nearly
  restored feasibility (peak $1.03$) but at a heavy visit cost ($G=-8.9\%$).
  \item \textbf{Still no free visits:} $G<0$ for $\Gamma=1$ on $4$ of $5$ cuts.
  \item \textbf{A silver lining for M3:} unlike the small instance, honest
  full-strength M3 was feasible at $\Gamma=1$ on all five cuts, so the planner's
  own deviation estimate is usable on this larger stream.
\end{itemize}

\plain{On a second, independent drifting instance the robust budget still trims
the worst station, but it is \emph{no longer clearly better} than simply
tightening the ceiling. And an exact solver confirms, rather than fixes, the
reliance on the heuristic. The paper's headline ``targeted protection dominates''
looks instance-specific, not a general law.}

% ======================================================================
\section{What it all means (plain summary)}
% ======================================================================

\begin{enumerate}
  \item \textbf{The audit is solid and important.} Reporting ``$0$ stations
  overloaded'' on late temporal cuts is misleading; that verdict mostly measures
  how little test demand there was. Feasibility should be judged on a common
  volume basis \eqref{eq:tnorm}.
  \item \textbf{The robust model is mathematically clean.} Bertsimas--Sim gives
  an exact, linear, easy-to-add protection for the workload constraint
  (Proposition~\ref{prop:1}), with a single intuitive dial $\Gamma$
  (Proposition~\ref{prop:2}). This part is textbook-correct.
  \item \textbf{But the practical case is narrow.}
  \begin{itemize}[nosep]
    \item It insures \emph{feasibility}, never visits ($G<0$ throughout).
    \item The honest ex-ante deviation model (M3) often under-sizes real drift,
    so on the small instance the ``win'' relied on a hand-tuned half-strength
    factor.
    \item Its strongest claim (targeted beats uniform tightening) did not
    replicate on a second drifting instance.
    \item Full protection does not fit the given capacity --- insuring big drift
    is a capacity-investment decision, not just a modelling one.
  \end{itemize}
  \item \textbf{Solver honesty.} At realistic scale, exact (CPLEX) and
  open-source (HiGHS) solvers agree exactly because both fall back on the
  heuristic incumbent; the gap comes from the weak linear relaxation of the visit
  objective, not the solver.
\end{enumerate}

\paragraph{One-sentence takeaway.}
The study correctly exposes a hidden flaw in how workload feasibility was
measured and offers a principled, correctly-implemented knob to protect against
demand drift --- but that knob buys feasibility (not visits), works only
partially inside the existing capacity, and is not universally better than the
simplest alternative, so its benefits should be stated as conditional rather than
general.

% ======================================================================
\section*{Reproducibility (where the code lives)}
% ======================================================================
\addcontentsline{toc}{section}{Reproducibility}

\begin{itemize}[nosep]
  \item Nominal/robust model \& solver (HiGHS): \texttt{src/milp\_highs\_robust.py}.
  \item CPLEX backend (twin model): \texttt{src/milp\_cplex\_robust.py}.
  \item Experiment harness (paired $\Gamma$-sweep, both backends via
  \texttt{-{}-backend}): \texttt{src/run\_bs\_robust\_experiment.py}.
  \item Stationary-benchmark adapter: \texttt{src/exp02a\_bs\_adapter.py}.
  \item Feasibility evaluator (as-run + normalized): \texttt{src/evaluate\_layout\_robust.py}.
  \item Results: \texttt{results/} (including
  \texttt{experiment\_iscf10kt\_cplex/} and \texttt{experiment\_iscf10kt\_highs/}).
\end{itemize}

\end{document}
